\section*{Chapter \ref{ch:rc:forward-rate-and-market-models} --- Forward-Rate and Market Models}

\begin{solution}{exo:rc:forward-rate-and-market-models:1}
$\alpha(t,T) = \sigma\int_t^T\sigma\,du = \sigma^2(T-t)$. It is the \omterm{def:rc:short-rate-models:hw}{Hull--White model} with
no mean reversion ($\kappa=0$), known as the Ho--Lee model: every forward moves by
the same normal shock.
\end{solution}

\begin{solution}{exo:rc:forward-rate-and-market-models:2}
Each forward is a martingale under its own payment measure with a deterministic
volatility, so it is lognormal there and its caplet is Black's formula. A swap
rate is a weighted sum of several forwards with stochastic weights; it is not
lognormal under its annuity measure, so Black is only an approximation.
\end{solution}

\begin{solution}{exo:rc:forward-rate-and-market-models:3}
$e^{-0.4\times8} = 0.0408$; with $\beta=0.02$, $e^{-0.16} = 0.852$.
\end{solution}

\begin{solution}{exo:rc:forward-rate-and-market-models:4}
At $t=0.5$ the next reset is $T_1$, so $m(t)=1$:
$\mu_2 = \sigma_2F_2\bigl(\frac{\delta\rho_{21}\sigma_1F_1}{1+\delta F_1}+\frac{\delta\sigma_2F_2}{1+\delta F_2}\bigr)$;
$F_1$ and $F_2$ appear ($F_0$ has already reset).
\end{solution}

\begin{solution}{exo:rc:forward-rate-and-market-models:5}
$\sigma_s^2 = \frac{1}{S^2}\sum w_iw_jF_iF_j\rho_{ij}\sigma^2 = 0.25\times0.04\times(2+2\rho)$:
20.00\% for $\rho=1$, 17.32\% for $\rho=0.5$.
\end{solution}

\begin{solution}{exo:rc:forward-rate-and-market-models:6}
A one-year swap starting in five years is a single annual forward, $F_5$: its
swaption is a caplet, priced by its own volatility, which both models calibrate
to the same target; correlation needs at least two forwards to matter.
\end{solution}

\begin{solution}{exo:rc:forward-rate-and-market-models:7}
EUR~2\,351\,823 with $\beta=0.02$, EUR~1\,840\,368 with $\beta=0.40$: a difference of
EUR~511\,456.
\end{solution}

\begin{solution}{exo:rc:forward-rate-and-market-models:8}
Each forward is a martingale only under its own payment measure; under a common
simulation measure all but one have a drift
(\cref{prop:rc:forward-rate-and-market-models:spotdrift}). With zero drifts the
simulation misprices every caplet except the one whose measure is the
numeraire's, by an error growing with maturity and volatility, and the
discounting by the spot account is inconsistent with the forwards.
\end{solution}

\begin{solution}{pb:rc:forward-rate-and-market-models:1}
\textbf{1.} 2.933\% and 4.102.
\textbf{2.} The integrated variance of each forward up to its reset (the $\phi_k$
given the volatility shape); not the correlation, and not how each forward's
variance is spread in time beyond what the chosen shape imposes.
\textbf{3.} 33.0\% and 23.1\%.
\textbf{4.} It depends on a single forward, $F_5$.
\textbf{5.} The shape parameters $a,b,c,d$, which move variance between early and late
times and so change swaptions of different expiries and forward volatilities.
\textbf{6.} 22.1\% and 17.3\%.
\textbf{7.} EUR~2.35 million and 1.84 million.
\textbf{8.} 22\%.
\textbf{9.} Within about two standard errors (under 0.9 volatility points on every
five-year-expiry swaption).
\textbf{10.} Something like chapter~6's Treasury estimate of 0.77 between two- and
ten-year changes: forward rates a few years apart are highly but not perfectly
correlated, which points to $\beta$ between the two extremes.
\textbf{11.} Swaptions on longer swaps (the \omterm{def:rc:sabr-in-rates-and-the-volatility-cube:matrix}{swaption matrix}), \omterm{def:rc:convexity-adjustments-and-constant-maturity-products:spread}{CMS spread options}
where traded.
\textbf{12.} About 0.16: with $\beta=0.16$ \omterm{def:rc:forward-rate-and-market-models:rebonato}{Rebonato's formula} gives 20\%.
\textbf{13.} The whole shape of the correlation matrix beyond one parameter, the
time dependence of correlation, and anything depending on joint moves at other
expiries.
\textbf{14.} It changes the swaption vegas' split across forwards and the hedge
ratios between caps and swaptions: a book hedged with caps against swaptions
carries correlation risk.
\textbf{15.} One factor moves all rates together: perfect correlation, the limit
$\beta\to0$.
\textbf{16.} The price at a $\beta$ calibrated to the quoted \omterm{def:rc:sabr-in-rates-and-the-volatility-cube:matrix}{swaption matrix}, with a
reserve covering a plausible range of correlation (here the difference between
$\beta$ values bracketing history and implied levels).
\textbf{17.} Track the calibrated and historical correlations daily, alert on jumps,
and report the P\&L from re-marking $\beta$ separately.
\textbf{18.} How it is chosen, what instruments it is calibrated to, its range, the
products sensitive to it and the reserve held.
\textbf{19.} Named result: \emph{the correlation you cannot see}: with the same caplets,
the five-into-five payer is worth EUR~2.35 million with $\beta=0.02$ and EUR~1.84
million with $\beta=0.40$, 22\% (EUR~511\,456) apart.
\textbf{20.} How rates of different maturities move together.
\end{solution}

\begin{solution}{iq:rc:forward-rate-and-market-models:1}
$\alpha(t,T)=\sigma_f(t,T)\int_t^T\sigma_f(t,u)\,du$ under the risk-neutral measure. It
comes from requiring every discounted bond price to be a martingale: the bond's
log-price drift, $r-\int\alpha+\frac12(\int\sigma_f)^2$, must equal $r$ for every
maturity; differentiate in $T$.
\iqlookfor{no arbitrage on bonds, and the derivation in two lines.}
\end{solution}

\begin{solution}{iq:rc:forward-rate-and-market-models:2}
Choose a measure (spot or terminal), write each forward's drift from the
covariance with the numeraire change, and step $\ln F_k$ with correlated normals,
a predictor--corrector drift and the numeraire rolled at resets. The drift is
there because each forward is a martingale only under its own payment measure.
\iqlookfor{common measure, drift from numeraire change, log stepping.}
\end{solution}

\begin{solution}{iq:rc:forward-rate-and-market-models:3}
A swap rate is an average of forwards with random weights; if forwards are
lognormal under their measures, a swap rate is not lognormal under its annuity
measure, and vice versa. One of the two sets is priced by Black only
approximately.
\iqlookfor{sums of lognormals are not lognormal.}
\end{solution}

\begin{solution}{iq:rc:forward-rate-and-market-models:4}
Caplets fix each forward's variance but not their correlation, which drives
swaption volatilities: the swaption prices are set by an arbitrary correlation
assumption, here up to 22\% apart. Calibrate to swaptions too.
\iqlookfor{correlation as the missing input.}
\end{solution}

\begin{solution}{iq:rc:forward-rate-and-market-models:5}
Hull--White: fast (trees), exact co-terminal fit with $\sigma(t)$, one factor so
perfectly correlated exercises, which understates the switch value. Market
model: flexible correlation and smile extensions, calibrated to caplets and
swaptions, but needs regression Monte Carlo for early exercise and is slower
and noisier. Use Hull--White in production, the market model as benchmark.
\iqlookfor{speed versus correlation, and a production/benchmark split.}
\end{solution}

\begin{solution}{iq:rc:forward-rate-and-market-models:6}
Discretisation bias of the drift: an Euler drift evaluated at the start of long
steps is wrong because it depends on forwards that move within the step. Use a
predictor--corrector, finer steps near resets, or a measure that puts the
numeraire's drift-free forward where it matters; check caplets against Black at
each step size.
\iqlookfor{drift discretisation and a convergence check.}
\end{solution}
